\section{Onde aprofundar}

\subsection*{Livros}
\small
\begin{longtable}{@{}p{1.8cm} p{4.4cm} p{2.6cm} p{4cm}@{}}
\toprule
Tema & Livro & Autor(es), ano & Por que ler \\
\midrule \endhead
Conduta & \citetitle{livro:martin2011} \pt{} \emph{O Codificador Limpo}
  & \citeauthor{livro:martin2011}, \citeyear{livro:martin2011}
  & Compromissos, estimativa, pressão e prática \\
Código & \citetitle{livro:martin2008} \pt{} \emph{Código Limpo}
  & \citeauthor{livro:martin2008}, \citeyear{livro:martin2008}
  & Nomes, funções, erros, classes e o catálogo de smells \\
Arquitetura & \citetitle{livro:martin2017} \pt{} \emph{Arquitetura Limpa}
  & \citeauthor{livro:martin2017}, \citeyear{livro:martin2017}
  & Dependency Rule e o framework como detalhe (caps.~22 e 32) \\
Java & \citetitle{livro:bloch2017}, 3ª ed. \pt{} \emph{Java Efetivo}
  & \citeauthor{livro:bloch2017}, \citeyear{livro:bloch2017}
  & Os itens da tabela deste capítulo \\
\bottomrule
\end{longtable}
\normalsize

\subsection*{Documentação}
\citetitle{doc:quarkus-cdi-reference}, seção \enquote{Native Executables
and Private Members}.

\subsection*{De quem fez o Java}
Nomes, bios e mais obras de cada pessoa no apêndice~\ref{ap:pessoas}.
\begin{itemize}
\obra*{video:bloch-api-design}{Design de API com exemplos bons e ruins do
  próprio JDK; vale para qualquer interface pública, inclusive portas da
  arquitetura hexagonal.}
\obra*{artigo:bumper-sticker-api}{As máximas da palestra em uma página.}
\end{itemize}

\subsection*{No IDE: blog da JetBrains}
Outros posts da JetBrains no apêndice~\ref{ap:jetbrains}.
\begin{itemize}
\obra*{artigo:jetbrains-code-smells-mutation}{Mutabilidade como code smell,
  e como records e imutabilidade resolvem.}
\obra*{artigo:jetbrains-code-smells-null}{O \texttt{null} como smell e as
  alternativas.}
\obra*{curso:jetbrains-guide-refactoring}{Refatorar para entender o
  código, com o IDE fazendo o trabalho mecânico.}
\end{itemize}

\subsection*{Pesquisa acadêmica}
Acesso (aberto ou pago) e links no apêndice~\ref{ap:pesquisa}.
\begin{itemize}
\pesquisa{paper:parnas-modules-1972}{Seis páginas que fundaram o
  information hiding: dividir o sistema pelas decisões que podem mudar.}
\pesquisa{paper:silva-why-we-refactor-2016}{Por que desenvolvedores
  refatoram, segundo eles mesmos, logo depois de refatorar.}
\pesquisa{paper:aniche-tese-usp-2016}{Tese do IME-USP: métricas e smells
  dependem do papel da classe (Controller, Repository).}
\pesquisa*{livro:valente2020}{Capítulo~5, sobre princípios de projeto,
  com leitura gratuita no site do livro.}
\pesquisa{doc:stanford-cs190}{O curso de design de software de Stanford que
  deu origem a \emph{A Philosophy of Software Design}.}
\end{itemize}
